\unit{Exposé 195 -- Technique de descente II}
\sid{195}
\src{369}
\seminar{12\textsuperscript{e} année, 1959/60, n\textsuperscript{o} 195\\Février 1960}
\paperhead{TECHNIQUE DE DESCENTE ET THÉORÈMES D'EXISTENCE EN GÉOMÉTRIE ALGÉBRIQUES\\
II. LE THÉORÈME D'EXISTENCE EN THÉORIE FORMELLE DES MODULES}%
  {par Alexander GROTHENDIECK}

\fsection{A}{Foncteurs représentables et pro-représentables}

\subsection*{1. Foncteurs représentables}
\addcontentsline{toc}{subsection}{1. Foncteurs représentables}

Soit $\mathcal C$ une catégorie. Pour tout $X\in\mathcal C$, soit $h_X$ le
foncteur contravariant de $\mathcal C$ dans la catégorie $(\operatorname{Ens})$
des ensembles,
\[
  h_X:\mathcal C\longrightarrow(\operatorname{Ens}),
\]
défini par la formule
\[
  h_X(Y)=\operatorname{Hom}(Y,X).
\]
Si on a un morphisme $X\longrightarrow X'$ dans $\mathcal C$, on en déduit de
façon évidente un homomorphisme $h_X\longrightarrow h_{X'}$ de foncteurs :
$h_X$ est un foncteur covariant en $X$, i.e. on a défini un foncteur covariant
canonique
\[
  h:\mathcal C\longrightarrow
  \operatorname{Hom}\bigl(\mathcal C^o,(\operatorname{Ens})\bigr)
\]
de $\mathcal C$ dans la catégorie des foncteurs covariants de la duale
$\mathcal C^o$ de $\mathcal C$ dans la catégorie des ensembles. Rappelons
alors :

\result{PROPOSITION}{1.1} Ce foncteur $h$ est \emph{pleinement fidèle}, en
d'autres termes, pour tout couple $X,X'$ d'objets de $\mathcal C$,
l'application naturelle
\[
  \operatorname{Hom}(X,X')\longrightarrow
  \operatorname{Hom}(h_X,h_{X'})
\]
est bijective.

En particulier, si un foncteur
$F\in\operatorname{Hom}(\mathcal C^o,(\operatorname{Ens}))$ est isomorphe à
un foncteur de la forme $h_X$, $X$ est déterminé à un isomorphisme unique près.
On dit alors que le foncteur $F$ est \emph{représentable}. La proposition
précédente signifie alors que le foncteur canonique $h$ définit une équivalence
de la catégorie $\mathcal C$ avec la sous-catégorie pleine de
$\operatorname{Hom}(\mathcal C^o,(\operatorname{Ens}))$ formée par les
foncteurs représentables. Ce fait est à la base de la notion de « solution
d'un Problème Universel », un tel problème consistant toujours à examiner si
un foncteur donné (contravariant comme ici, ou covariant dans le cas dual) de
$\mathcal C$ dans $(\operatorname{Ens})$ est représentable.

\src{370}
Notons de plus que, par définition même de la notion de produit dans une
catégorie [1], le foncteur $h:X\longmapsto h_X$ commute aux produits chaque
fois qu'ils existent (et plus généralement aux limites projectives finies ou
infinies, en particulier aux produits fibrés, à la formation de « noyaux » [2],
etc., chaque fois qu'ils existent) : on a un isomorphisme de foncteurs
\[
  h_{X\times X'}\xrightarrow{\sim}h_X\times h_{X'}
\]
chaque fois que $X\times X'$ existe, i.e. on a des bijections fonctorielles en
$Y$
\[
  h_{X\times X'}(Y)\xrightarrow{\sim}h_X(Y)\times h_{X'}(Y).
\]
En particulier, la donnée d'un morphisme
\[
  X\times X'\longrightarrow X''
\]
dans $\mathcal C$ (i.e. d'une « loi de composition » dans $\mathcal C$ entre
$X,X',X''$) est équivalente à la donnée d'un morphisme
\[
  h_{X\times X'}=h_X\times h_{X'}\longrightarrow h_{X''},
\]
i.e. à la donnée pour tout $Y\in\mathcal C$ d'une loi de composition
d'ensembles
\[
  h_X(Y)\times h_{X'}(Y)\longrightarrow h_{X''}(Y)
\]
de telle façon que, pour tout morphisme $Y\longrightarrow Y'$ dans
$\mathcal C$, le système des applications ensemblistes
\[
  h_{X^{(i)}}(Y)\longleftarrow h_{X^{(i)}}(Y')
  \qquad(\text{pour }i=0,1,2)
\]
soit un morphisme pour les deux lois de composition, relatives à $Y$ et $Y'$.
On voit de cette façon que la notion de structure de « $\mathcal C$-groupe »,
« $\mathcal C$-anneau », etc. sur un objet $X$ de $\mathcal C$ s'exprime de la
façon la plus commode (en théorie comme en pratique) en disant que pour tout
$Y\in\mathcal C$, on a une loi de groupes (resp. anneau, etc.) au sens usuel
sur l'ensemble $h_X(Y)$, les applications
$h_X(Y)\longleftarrow h_X(Y')$ correspondant à des morphismes
$Y\longrightarrow Y'$ devant être des homomorphismes de groupes (resp.
anneaux, etc.). C'est la façon la plus intuitive et la plus commode, par
exemple, pour définir les divers groupes classiques $G_a$, $G_m$,
$\operatorname{Gl}(n)$, etc. sur un pré-schéma $S$ de base quelconque et pour
écrire les relations classiques entre ces groupes, pour mettre sur le schéma
affine $V(\mathcal F)$ au-dessus de $S$, défini par un faisceau quasi-cohérent
$\mathcal F$, une structure de « fibré vectoriel », pour définir et étudier les
diverses variétés de drapeaux (grassmanniennes, fibrés projectifs) associés,
etc. ; le yoga général étant d'identifier purement et simplement, à l'aide du
foncteur canonique $h$, les objets de $\mathcal C$ à des foncteurs
contravariants particuliers, les foncteurs représentables,

\src{371}
de $\mathcal C$ dans la catégorie des ensembles.

Le procédé habituel de renversement de flèches, qui s'impose par exemple, dans
le cas des schémas affines pour passer du langage géométrique au langage de
l'algèbre commutative, conduit à dualiser les considérations précédentes et en
particulier, à introduire aussi les foncteurs covariants
$\mathcal C\longrightarrow(\operatorname{Ens})$ représentables, i.e. de la
forme
\[
  Y\longmapsto\operatorname{Hom}(X,Y)=h'_X(Y).
\]

\subsection*{2. Foncteurs pro-représentables, pro-objets}
\addcontentsline{toc}{subsection}{2. Foncteurs pro-représentables, pro-objets}

Soit $\widetilde X=(X_i)_{i\in I}$ un système projectif d'objets de
$\mathcal C$ ; il lui correspond un foncteur covariant
\[
  h'_{\widetilde X}=\varinjlim_i h'_{X_i},
\]
de façon plus explicite
\[
  h'_{\widetilde X}(Y)=\varinjlim_i h'_{X_i}(Y)
  =\varinjlim_i\operatorname{Hom}(X_i,Y),
\]
de $\mathcal C$ dans $(\operatorname{Ens})$. Un foncteur de $\mathcal C$ dans
$(\operatorname{Ens})$ qui est isomorphe à un foncteur de ce type avec $I$
filtrant est dit \emph{pro-représentable}. D'après le numéro précédent, ce
sont exactement les foncteurs isomorphes à des limites inductives filtrantes
de foncteurs représentables. Soit
$\widetilde X'=(X'_j)_{j\in J}$ un deuxième système projectif filtrant dans
$\mathcal C$ (construit sur un deuxième ensemble d'indices préordonné filtrant
$J$) ; on vérifie aisément qu'on a une bijection canonique (généralisant la
proposition 1.1)
\[
  \operatorname{Hom}(h'_{\widetilde X'},h'_{\widetilde X})
  =\varprojlim_j\varinjlim_i\operatorname{Hom}(X_i,X'_j).
\]
Cela amène à introduire la catégorie $\operatorname{Pro}(\mathcal C)$ des
pro-objets de $\mathcal C$ ; ses objets sont les systèmes projectifs d'objets
de $\mathcal C$ (sur des ensembles d'indices préordonnés filtrants variables),
et si $\widetilde X=(X_i)_{i\in I}$ et
$\widetilde X'=(X'_j)_{j\in J}$ sont deux tels objets, on pose
\[
  \operatorname{ProHom}(\widetilde X,\widetilde X')
  =\varprojlim_j\varinjlim_i\operatorname{Hom}(X_i,X'_j),
\]
la composition des pro-homomorphismes étant évidente. Par construction même,
$\widetilde X\longmapsto h'_{\widetilde X}$ peut être considéré comme un
foncteur contravariant en $\widetilde X$, établissant une équivalence de la
catégorie duale de la catégorie $\operatorname{Pro}(\mathcal C)$ des
pro-objets de $\mathcal C$, avec la catégorie des foncteurs covariants
pro-représentables de $\mathcal C$ dans $(\operatorname{Ens})$. Bien entendu,
un objet $X$ de $\mathcal C$ définit canoniquement un pro-objet, dénoté

\src{372}
encore par $X$, de sorte que $\mathcal C$ est équivalente à une sous-catégorie
pleine de $\operatorname{Pro}(\mathcal C)$. Si alors
$\widetilde X=(X_i)_{i\in I}$ est un pro-objet quelconque de $\mathcal C$, on
aura (avec l'identification précédente)
\[
  \widetilde X=\varprojlim_i X_i,
\]
la limite projective étant prise dans $\operatorname{Pro}(\mathcal C)$
(puisque $h'_{\widetilde X}=\varinjlim_i h'_{X_i}$). On fera attention que,
lors même que la limite projective des $X_i$ existe dans $\mathcal C$, elle
sera généralement non isomorphe à la limite projective $\widetilde X$ dans
$\operatorname{Pro}(\mathcal C)$, comme il est déjà évident dans le cas où
$\mathcal C$ est la catégorie des ensembles. On notera que par définition
même, $\varprojlim_{\mathcal C}X_i=L$ est défini par la condition que le
foncteur en $Y\in\mathcal C$
\[
  \varprojlim_i\operatorname{Hom}_{\mathcal C}(Y,X_i)
  =\operatorname{Hom}_{\operatorname{Pro}(\mathcal C)}(Y,\widetilde X)
\]
et à valeurs dans $(\operatorname{Ens})$ soit représentable à l'aide de $L$,
i.e. soit isomorphe à $\operatorname{Hom}_{\mathcal C}(Y,L)$ ; par suite,
$\varprojlim_{\mathcal C}X_i$ est défini déjà en termes du pro-objet
$\widetilde X$, et de façon précise dépend fonctoriellement du pro-objet
$\widetilde X$ chaque fois qu'elle est définie ; il n'y a donc pas
d'inconvénient à la dénoter par $\varprojlim_{\mathcal C}(\widetilde X)$. Si
les limites projectives dans $\mathcal C$ existent toujours,
$\varprojlim_{\mathcal C}(\widetilde X)$ est un foncteur de
$\operatorname{Pro}(\mathcal C)$ dans $\mathcal C$, et il y a un
homomorphisme canonique de foncteurs
$\varprojlim_{\mathcal C}(\widetilde X)\longrightarrow\widetilde X$.

Comme tout foncteur, covariant pour fixer les idées,
\[
  F:\mathcal C\longrightarrow\mathcal C'
\]
se prolonge de façon évidente en un foncteur
\[
  \operatorname{Pro}(F):\operatorname{Pro}(\mathcal C)
  \longrightarrow\operatorname{Pro}(\mathcal C'),
\]
il s'ensuit que, si dans $\mathcal C'$ les limites projectives existent
toujours, $F$ définit aussi canoniquement un foncteur composé
\[
  \overline F:\operatorname{Pro}(\mathcal C)\longrightarrow\mathcal C',
  \qquad
  \overline F=\varprojlim_{\mathcal C'}\circ\operatorname{Pro}(F),
\]
transformant $\widetilde X=(X_i)_{i\in I}$ en
$\varprojlim_{\mathcal C'}F(X_i)$.

Un pro-objet $\widetilde X$ est dit \emph{pro-objet strict} s'il est isomorphe
à un pro-objet $(X_i)_{i\in I}$ où les morphismes de transition
$X_i\longrightarrow X_j$ sont des épimorphismes ; un foncteur défini par un
tel objet est dit \emph{strictement pro-représentable}. On peut alors exiger de
plus que $I$ soit ordonné filtrant, et que tout épimorphisme
$X_i\longrightarrow X'$ soit équivalent à un épimorphisme

\src{373}
$X_i\longrightarrow X_j$ pour $j\in I$ convenable (déterminé de façon unique
par cette condition). Sous ces conditions, le système projectif
$(X_i)_{i\in I}$ est déterminé à un isomorphisme unique près (au sens usuel
d'isomorphismes de systèmes projectifs). Il en résulte que dans
$\operatorname{Pro}(\mathcal C)$ la limite projective d'un système projectif de
pro-objets stricts $\widetilde X^{(\alpha)}$ existe toujours et qu'avec les
notations ci-dessus pour $F$, $\overline F$, on aura
\[
  \overline F\left(\varprojlim_\alpha\widetilde X^{(\alpha)}\right)
  =\varprojlim_{\alpha,\mathcal C'}
  \overline F\left(\widetilde X^{(\alpha)}\right).
\]
En particulier, si tout pro-objet de $\mathcal C$ est strict (cf. le numéro
suivant), le foncteur prolongé $\overline F$ commute aux limites projectives.

\subsection*{3. Caractérisation des foncteurs pro-représentables}
\addcontentsline{toc}{subsection}{3. Caractérisation des foncteurs pro-représentables}

Soient $\mathcal C$, $\mathcal C'$ deux catégories où les limites projectives
finies (i.e. relatives à des ensembles préordonnés finis, non nécessairement
filtrants) existent, ou ce qui revient au même, où les produits finis et les
produits fibrés finis existent (ce qui implique en particulier l'existence
d'un « objet unité à droite » $e$, tel que $\operatorname{Hom}(X,e)$ soit
réduit à un élément pour tout $X$). Soit $F$ un foncteur covariant de
$\mathcal C$ dans $\mathcal C'$. Les conditions suivantes sont équivalentes :

\begin{enumerate}
\item[(i)] $F$ permute aux limites projectives finies ;
\item[(ii)] $F$ permute aux produits finis et aux produits fibrés finis ;
\item[(iii)] $F$ permute aux produits finis, et pour tout diagramme exact
\[
  X\longrightarrow X'\rightrightarrows X''
\]
dans $\mathcal C$ ([3], A, définition 2.1), le diagramme transformé par $F$
\[
  F(X)\longrightarrow F(X')\rightrightarrows F(X'')
\]
est exact.
\end{enumerate}

On dit alors que $F$ est \emph{exact à gauche}.

Par la suite, on suppose que dans $\mathcal C$ les limites projectives finies
existent. Il est alors immédiat sur les définitions qu'un foncteur
représentable est exact à gauche, et par passage à la limite qu'un foncteur
pro-représentable est exact à gauche.

Pour obtenir une réciproque, soit
\[
  F:\mathcal C\longrightarrow(\operatorname{Ens})
\]
un foncteur covariant, soient $X\in\mathcal C$ et $\xi\in F(X)$. On dit que
$\xi$ (ou le couple \src{374}$(X,\xi)$) est \emph{minimal} si pour tout
$X'\in\mathcal C$ et $\xi'\in F(X')$, et tout monomorphisme strict ([3], A, 2)
$u:X'\longrightarrow X$ tel que $\xi=F(u)(\xi')$, $u$ est un isomorphisme. On
dit d'autre part qu'un couple $(X,\xi)$ \emph{domine} $(X'',\xi'')$
($\xi\in F(X)$, $\xi''\in F(X'')$) s'il existe un morphisme
$v:X\longrightarrow X''$ tel que $\xi''=F(v)(\xi)$ ; si $\xi$ est minimal et
si $F$ est exact à gauche, ce morphisme $v$ est alors unique ; si $\xi''$ est
minimal, $v$ est surjectif. On en déduit facilement la proposition suivante :

\result{PROPOSITION}{3.1} Pour que $F$ soit strictement pro-représentable, il
faut et il suffit qu'il satisfasse aux deux conditions :
\begin{enumerate}
\item[(i)] $F$ est exact à gauche ;
\item[(ii)] tout couple $(X,\xi)$, avec $\xi\in F(X)$, est dominé par un couple
minimal.
\end{enumerate}

Cette dernière condition est vide si tout objet de $\mathcal C$ est artinien
(en prenant un sous-objet $X'$ de $X$ minimal parmi ceux pour lesquels il
existe $\xi'\in F(X')$ tel que $\xi$ soit l'image de $\xi'$). D'où :

\result{COROLLAIRE}{} Soit $\mathcal C$ une catégorie dont les objets sont
artiniens et où les limites projectives finies existent. Alors les foncteurs
pro-représentables de $\mathcal C$ dans $(\operatorname{Ens})$ sont exactement
les foncteurs exacts à gauche, et ils sont même strictement
pro-représentables.

Ce dernier fait signifie aussi que \emph{tout pro-objet de $\mathcal C$ est
strict}.

\subsection*{4. Exemple : groupes du type galoisien, groupes pro-algébriques}
\addcontentsline{toc}{subsection}{4. Groupes du type galoisien, groupes pro-algébriques}

Si $\mathcal C$ est la catégorie des groupes finis ordinaires,
$\operatorname{Pro}(\mathcal C)$ est équivalente à la catégorie des groupes
topologiques compacts totalement discontinus. Ce sont des groupes de ce type
et leurs généralisations, obtenues en remplaçant les groupes finis ordinaires
par des schémas en groupes finis au-dessus d'un pré-schéma de base donné (par
exemple, les groupes algébriques finis sur un corps $k$), qui tiendront lieu de
groupes fondamentaux, d'homotopie, d'homologie absolus et relatifs pour les
pré-schémas. Dans tous ces exemples, le corollaire à la proposition 3.1
s'applique, et c'est bien par le foncteur associé que le $\pi_1$ doit se
définir [2]. Il en est de même en partant de la catégorie des groupes
algébriques ou quasi-algébriques sur un corps (ou plus généralement sur un
pré-schéma noethérien) : on trouve les « groupes pro-algébriques » de Serre
[4].

\subsection*{5. Exemple : « variétés formelles »}
\addcontentsline{toc}{subsection}{5. Exemple : variétés formelles}

Soient $\Lambda$ un anneau noethérien, $\mathcal C$ la catégorie des
$\Lambda$-algèbres $A$ qui sont des modules de type fini artiniens sur
$\Lambda$ (ou plus brièvement,

\src{375}
des $\Lambda$-algèbres artiniennes). On est sous les conditions du corollaire
à la proposition 3.1. La catégorie $\operatorname{Pro}(\mathcal C)$ est ici
équivalente à la catégorie des algèbres topologiques $O$ sur $\Lambda$
isomorphes à des limites projectives topologiques
\[
  O=\varprojlim_i O_i
\]
d'algèbres $O_i\in\mathcal C$, i.e. dont la topologie est linéaire, séparée et
complète, et telle que pour tout idéal ouvert $J_i$ de $O$, $O/J_i$ soit une
algèbre artinienne sur $\Lambda$. Le foncteur
$\mathcal C\longrightarrow(\operatorname{Ens})$ associé à une telle algèbre
n'est autre que
\begin{align*}
  F(A)=h'_O(A)
  &=\{\text{homomorphismes continus de $\Lambda$-algèbres topologiques }
       O\longrightarrow A\}\\
  &=\varinjlim_i
    \operatorname{Hom}_{\Lambda\text{-algèbres}}(O_i,A).
\end{align*}

On notera d'ailleurs que la catégorie $\mathcal C$ est essentiellement le
produit des catégories analogues, correspondant aux anneaux locaux complétés
des $\Lambda_{\mathfrak m}$ pour les idéaux maximaux $\mathfrak m$ de
$\Lambda$ ; on peut donc si on le désire se limiter au cas où $\Lambda$ est un
tel anneau local complet. En tous cas, $O$ se décompose canoniquement en le
produit topologique de composants locaux, correspondant aux « points » du
schéma formel [2] défini par $O$. Un tel point est défini par un objet $\xi$
d'un $F(K)$, où $K\in\mathcal C$ est un corps (par exemple, le corps résiduel
du composant local envisagé), deux couples $(\xi,K)$ et $(\xi',K')$
définissant le même point si et seulement s'ils sont dominés par un même
$(\xi'',K'')$, ou encore s'ils dominent un même $(\xi''',K''')$. (Si les
$\Lambda/\mathfrak m$ sont algébriquement clos, il suffit de prendre l'ensemble
somme des $F(\Lambda/\mathfrak m)$.)

Il importe de donner des conditions pour que le composant local $O_\xi$ de
$O$ correspondant à un $\xi\in F(K)$ soit un anneau noethérien. Lorsque
$\Lambda$ est un anneau local complet (noethérien, on le rappelle), il revient
au même de dire que $O_\xi$ est isomorphe à un anneau quotient d'un anneau de
séries formelles $\Lambda[[t_1,\ldots,t_n]]$ sur $\Lambda$. Pour donner un tel
critère, introduisons (pour tout anneau $A$) la $A$-algèbre $I_A$ des « nombres
duaux » de $A$ par
\[
  I_A=A[t]/t^2A[t].
\]
Soit $\varepsilon:I_A\longrightarrow A$ l'homomorphisme d'augmentation ; il
définit (si $A\in\mathcal C$) une application
\[
  F(\varepsilon):F(I_A)\longrightarrow F(A),
\]

\src{376}
et utilisant le fait que $F$ est exact à gauche, on définit de façon
intrinsèque une structure de $A$-module dans la partie
\[
  F(I_A,\xi)=F(\varepsilon)^{-1}(\xi)
\]
de $F(I_A)$ formée des $\xi'\in F(I_A)$ qui « se réduisent suivant $\xi$ » ;
utilisant la forme explicite de $F$ en termes de $O$, on trouve que ce
$A$-module s'identifie à
\[
  \operatorname{Hom}_{\Lambda}(\mathfrak m_\xi/\mathfrak m_\xi^2,A),
\]
où $\mathfrak m_\xi$ est le noyau de l'homomorphisme
$\xi:O\longrightarrow A$, i.e. si $A$ est un corps, l'idéal maximal du
composant local $O_\xi$ de $O$. On en déduit aussitôt la proposition suivante :

\result{PROPOSITION}{5.1} Soit $\xi\in F(K)$, où $K\in\mathcal C$ est un
corps. Pour que le composant local $O_\xi$ correspondant de $O$ soit un anneau
noethérien, il faut et il suffit que l'ensemble $F(I_K,\xi)$ des éléments de
$F(I_K)$ se réduisant suivant $\xi$ soit un espace vectoriel de dimension
finie sur $K$. Sous ces conditions, on a un isomorphisme canonique
\[
  F(I_K,\xi)=
  \operatorname{Hom}\!\left(
    \frac{\mathfrak m_\xi/\mathfrak m_\xi^2}
         {\operatorname{im}(\mathfrak n_\xi/\mathfrak n_\xi^2)},K
  \right),
\]
où $\mathfrak n_\xi$ est l'idéal maximal de $\Lambda$, noyau de
l'homomorphisme $\Lambda\longrightarrow K$.
\footnote{Le texte imprimé donne ici une formule qui n'est correcte que
lorsque $\Lambda$ est un corps. L'erratum de 1962 prescrit, dans le cas
général, le quotient de $\mathfrak m_\xi/\mathfrak m_\xi^2$ par l'image de
$\mathfrak n_\xi/\mathfrak n_\xi^2$.}
En particulier, la dimension du $K$-espace vectoriel $F(I_K,\xi)$ est égale à
la dimension de l'espace vectoriel ainsi corrigé sur le corps
$O_\xi/\mathfrak m_\xi=K(\xi)$.

Supposons que $O_\xi$ soit noethérien, et supposons pour simplifier l'écriture
que $\Lambda$ soit local complet, et $O=O_\xi$. On dit que $O$ est
\emph{simple sur $\Lambda$} si $O$ est une algèbre finie et étale sur
l'algèbre complétée du localisé de $\Lambda[t_1,\ldots,t_n]$ en un idéal
maximal de cet anneau, induisant l'idéal maximal de $\Lambda$ ; lorsque
l'extension résiduelle de $O$ sur $\Lambda$ est triviale (par exemple, le
corps résiduel de $\Lambda$ est algébriquement clos), cela équivaut à dire que
$O$ lui-même est isomorphe à une telle algèbre de séries formelles.
\footnote{L'erratum de 1962 précise que cette définition de « simple sur
$\Lambda$ » n'est correcte que lorsque l'extension résiduelle $k'/k$ est
séparable ; dans le cas général, il renvoie à SGA III, 1.1.}
Enfin, si on ne suppose plus nécessairement $O$ noethérien, on dira encore que
$O$ est simple sur $\Lambda$ si $O$ est isomorphe à une limite projective
topologique de $\Lambda$-algèbres quotients qui sont noethériennes et
$\Lambda$-simples au sens précédent. On généralise aussitôt au cas où
$\Lambda$, $O$ ne sont plus supposés locaux. Ceci dit, on a la proposition
suivante :

\result{PROPOSITION}{5.2} Pour que $O$ soit simple sur $\Lambda$, il faut et il
suffit que le foncteur $F$ associé transforme épimorphismes en épimorphismes.

Cela signifie donc que pour tout homomorphisme surjectif
$A\longrightarrow A'$ dans $\mathcal C$,
$F(A)\longrightarrow F(A')$ est surjectif. Bien entendu, il suffit de vérifier
cette condition quand $A$ est local, et (en procédant de proche en proche)
quand l'idéal de $A$ noyau de $A\longrightarrow A'$ est annulé par l'idéal
maximal de $A$. Cela ramène, en pratique, à vérifier qu'une certaine
obstruction, liée à des invariants infinitésimaux

\src{377}
de la situation donnant naissance au foncteur $F$, est nulle, problème qui est
de nature cohomologique.

Pour finir, disons quelques mots, dans le contexte précédent, des \emph{anneaux
de définition}. Soit toujours $F$ un foncteur de $\mathcal C$ dans
$(\operatorname{Ens})$, pro-représentable à l'aide d'une $\Lambda$-algèbre
topologique $O$. Alors pour tout $A\in\mathcal C$ et tout $\xi\in F(A)$, il
existe un plus petit sous-anneau $A'$ de $A$ tel que $\xi$ soit image d'un
élément $\xi'$ de $F(A')$ (lequel est alors uniquement déterminé) : en effet,
il suffit d'interpréter $\xi$ au moyen d'un homomorphisme de $O$ dans $A$, et
de prendre pour $A'$ l'image de $O$ par ce dernier. On dit alors que $A'$ est
l'anneau de définition de l'objet $\xi\in F(A)$. Si
$u:A\longrightarrow B$ est un homomorphisme d'algèbres et si
$\eta=F(u)(\xi)$, alors l'anneau de définition de $\eta$ est l'image par $u$ de
l'anneau de définition de $\xi$. Lorsqu'on part d'un foncteur $F$ de
$\mathcal C$ dans $(\operatorname{Ens})$, l'existence des anneaux de définition
et les propriétés qu'on vient d'en donner sont à peu de choses près
équivalentes au fait que $F$ soit pro-représentable ; c'est dire qu'elles sont
le plus souvent loin d'être triviales.

\fsection{B}{Les deux théorèmes d'existence}

Gardons les notations de A, paragraphe 5, et partons d'un foncteur covariant
\[
  F:\mathcal C\longrightarrow(\operatorname{Ens}) ;
\]
on cherche des critères maniables pour que $F$ soit pro-représentable, i.e.
exprimable à l'aide d'une $\Lambda$-algèbre $O$ comme plus haut. En vertu du
corollaire de A, proposition 3.1, il faut et il suffit pour cela que $F$ soit
exact à gauche. Dans l'état actuel de la technique de descente (cf. questions
posées dans [3], p. 9), ce critère n'est pas directement vérifiable sous cette
forme dans les cas les plus importants, et on a besoin de critères en apparence
moins exigeants.

\result{THÉORÈME}{1} Pour que le foncteur $F$ soit pro-représentable, il faut
et il suffit qu'il satisfasse les deux conditions suivantes :
\begin{enumerate}
\item[(i)] $F$ permute aux produits finis ;
\item[(ii)] pour toute algèbre $A\in\mathcal C$ et tout homomorphisme
$A\longrightarrow A'$ dans $\mathcal C$ tel que le diagramme
\[
  A\longrightarrow A'\rightrightarrows A'\otimes_A A'
\]
soit exact ([3], A, définition 1.2), le diagramme transformé
\[
  F(A)\longrightarrow F(A')\rightrightarrows F(A'\otimes_A A')
\]

\src{378}
est exact.
\end{enumerate}
De plus, il suffit de vérifier (ii) quand $A$ est local et quand de plus on est
dans l'un des deux cas suivants :
\begin{enumerate}
\item[(a)] $A'$ est un module libre sur $A$ ;
\item[(b)] le module quotient $A'/A$ est un $A$-module de longueur 1.
\end{enumerate}

La démonstration de ce théorème est assez délicate et ne peut être esquissée
ici. Bornons-nous à signaler qu'elle repose essentiellement sur une étude des
relations d'équivalence (au sens des catégories) dans le spectre d'une algèbre
artinienne (étude qui pose encore plusieurs problèmes, dont la solution semble
indispensable pour le développement ultérieur de la théorie).

Dans les applications, la vérification de la condition (i) est toujours
triviale. Celle de (ii) se décompose en deux : le cas où $A'$ est un
$A$-module libre relève de la technique de descente par morphismes plats ([3],
théorèmes 1, 2 et 3) et n'offre pas de difficulté. Pour traiter le cas (b), on
utilisera le résultat suivant :

\result{THÉORÈME}{2} Soient $A$ un anneau local artinien d'idéal maximal
$\mathfrak m$, $A'$ une $A$-algèbre contenant $A$ telle que
$\mathfrak mA'\subset A$ et que
\[
  A\longrightarrow A'\rightrightarrows A'\otimes_A A'
\]
soit exact (ce qui est le cas en particulier si $A'/A$ est un $A$-module de
longueur 1). Soit $\mathcal F$ la catégorie fibrée ([3], A, définition 1.1)
des faisceaux quasi-cohérents et plats sur des pré-schémas variables. Alors le
morphisme $\operatorname{Spec}(A')\longrightarrow\operatorname{Spec}(A)$ est
un morphisme de $\mathcal F$-descente strict ([3], A, définition 1.7).

En d'autres termes, la donnée d'un $A$-module plat $M$ est complètement
équivalente à la donnée d'un $A'$-module plat $M'$, muni d'un
$A'\otimes_AA'$-isomorphisme de $M'\otimes_AA'$ sur $A'\otimes_AM'$
satisfaisant la condition de transitivité habituelle pour une donnée de
descente (loco citato).

Le théorème 2 se démontre en prouvant d'abord que
\[
  H^i(A'/A,G_a)=0\qquad\text{pour }i\geq 1
\]
([3], A, 4(e)), l'hypothèse $\mathfrak mA'\subset A$ permettant de se ramener
aisément au cas où $A$ est un corps $k$ (savoir $A/\mathfrak m$). On applique
alors les équivalences signalées dans [3], page 16.

\fsection{C}{Applications à quelques cas particuliers}

\subsection*{1. Remarques générales sur les foncteurs représentés par des pré-schémas}
\addcontentsline{toc}{subsection}{1. Remarques générales sur les foncteurs représentés par des pré-schémas}

Soit $S$ un pré-schéma localement noethérien. Un pré-schéma $X$ sur $S$ est dit
\src{379}localement de type fini sur $S$ si, pour tout $x\in X$ se projetant en
$y\in S$,\footnote{Le texte imprimé porte ici « $y\in Y$ » et emploie de
nouveau $Y$ à la fin du paragraphe. Le morphisme considéré étant
$X\longrightarrow S$, on lit $S$ dans les deux occurrences.} il existe un
voisinage affine de $y$, d'anneau $A$, et un voisinage affine de $x$ au-dessus
de ce dernier, d'anneau $B$, tels que $B$ soit une $A$-algèbre de type fini.
On trouve de nombreux exemples importants de pré-schémas localement de type
fini sur $S$, qui ne sont pas de type fini sur $S$, comme solutions de
problèmes universels classiques ; ainsi il importe de pouvoir considérer le
schéma de Picard d'une courbe comme réunion d'une infinité de composantes
connexes (qu'il faut se garder de confondre avec la composante connexe de
l'élément neutre, i.e. la « variété de Picard »). Il est alors parfois commode
de se placer dans la catégorie $\mathcal C$ des pré-schémas localement de type
fini sur $S$, pour y examiner la question de la représentabilité d'un foncteur
contravariant $F$. Le but principal de ces exposés est de développer une
technique générale permettant de reconnaître si un tel foncteur $F$ est
représentable, et d'étudier les propriétés du $S$-pré-schéma $X$ correspondant
à l'aide de celles de $F$. Signalons en passant que, dans cette étude, on
trouve des exemples non pathologiques de pré-schémas sur $S$ non séparés sur
$S$, notamment comme « pré-schémas de Picard » d'excellents $S$-schémas ; il
faut donc se garder de bannir de la géométrie algébrique les pré-schémas qui ne
sont pas des schémas.

Soit $X$ un pré-schéma localement de type fini sur $S$, et soit
\[
  F(Y)=\operatorname{Hom}_S(Y,X)
\]
le foncteur contravariant associé. On peut considérer la restriction $F_0$ de
$F$ à la sous-catégorie $\mathcal C_0$ de $\mathcal C$ formée des
pré-schémas $Y$ sur $S$ qui sont artiniens et finis sur $S$ : lorsque
$S=\operatorname{Spec}(\Lambda)$, $\mathcal C_0$ est donc la catégorie duale
de la catégorie des $\Lambda$-algèbres artiniennes considérée dans B. Si
$Y=\operatorname{Spec}(A)$, où $A$ est un anneau local artinien, $Y$ est
réduit à un seul point $y$ placé au-dessus d'un point fermé $s$ de $S$, et un
$S$-homomorphisme de $Y$ dans $X$ (i.e. un élément de $F(Y)$) est défini par
la donnée d'un point $x\in X$ au-dessus de $s$, et d'un
$\mathcal O_{s,S}$-homomorphisme de $\mathcal O_{x,X}$ dans $A$. S'il existe
un tel homomorphisme, alors $x$ est nécessairement un point fermé de $X$ (son
corps résiduel étant algébrique sur celui de $s$). Cela montre donc que la
restriction $F_0$ de $F$ aux « $\Lambda$-algèbres artiniennes » est
pro-représentable, et est représentée par la $\Lambda$-algèbre topologique
dont les composants locaux sont les complétés
$\widehat{\mathcal O}_{x,X}$ des anneaux locaux de $X$ aux points $x$ de $X$
qui sont fermés et se trouvent au-dessus de points fermés de $S$. Cela montre
que la seule connaissance de $F_0$ donne des renseignements précis sur la
structure de $X$ (savoir la structure des complétés de ses anneaux locaux aux
points indiqués). On notera que, même dans le cas où $S$

\src{380}
est le spectre d'un corps algébriquement clos, ce n'est que grâce à la
considération systématique de « variétés » $Y$ telles que
$\mathcal O_Y$ peut admettre des éléments nilpotents, en particulier en
travaillant avec les spectres d'anneaux artiniens locaux, qu'on peut arriver à
la « bonne formulation » des problèmes universels classiques, et à en saisir
l'aspect « infinitésimal ».

Quand on part d'un foncteur $F$ donné à l'avance, dont on veut décider s'il est
représentable, l'étude du foncteur $F_0$ (à l'aide des théorèmes 1 et 2)
donnera des indications quasi complètes ; soit, comme il arrive fréquemment
(en testant simplement, par exemple, la nature des ensembles
$F(I_K,\xi)$ et leur comportement fonctoriel, cf. A), qu'on constate que
$F_0$ déjà n'est pas pro-représentable (ce qui explique l'échec des tentatives
faites jusqu'à présent pour définir de façon plus ou moins naturelle des
variétés de modules pour la classification des fibrés vectoriels de rang
$>1$) ; soit qu'on arrive à vérifier que $F_0$ est bien représentable, mais
que les espaces vectoriels $F(I_K,\xi)$ ne sont pas de dimension finie,
auquel cas il faut se contenter de la solution « formelle » ; soit enfin que
$F_0$ est bien représentable par un produit d'anneaux locaux complets
noethériens, ce qui donne de très fortes présomptions pour que $F$ soit
lui-même représentable et, joint à des propriétés analogues mais de nature
plus globale que nous pensons développer par la suite, suffira sans doute à
entraîner qu'il en est effectivement ainsi. Enfin, on rencontre des problèmes
géométriques intéressants (voir paragraphes 4 et 5 ci-dessous) où on ne dispose
que du foncteur $F_0$ (ne provenant d'aucun foncteur « global » $F$), et où on
s'estimera heureux quand on aura pu lui associer un « schéma formel de
modules ».

Pour terminer ces généralités, indiquons comment la théorie des schémas
explique des anomalies apparentes, telle la surface de Igusa $V$ dont la
« variété de Picard » $P$ est réduite à un point, et pour laquelle pourtant
on a
\[
  H^1(V,\mathcal O_V)\ne 0 ;
\]
dans ce cas, $P$ est un groupe « purement infinitésimal » non réduit au groupe
unité, i.e. défini par une algèbre locale $O$ de rang fini sur le corps de base
$k$ et muni d'une application diagonale correspondant à la structure
multiplicative de $P$ ; si $\mathfrak m$ est l'idéal maximal de $O$, le dual
de $\mathfrak m/\mathfrak m^2$ est canoniquement isomorphe à
$H^1(V,\mathcal O_V)$ (cf. le paragraphe 3 ci-dessous). Ce n'est que lorsque le
groupe de Picard est un groupe algébrique au sens classique, i.e. simple sur le
corps de base $k$, que la dimension de $H^1(V,\mathcal O_V)$ (toujours égale à
celle de $\mathfrak m/\mathfrak m^2$) est égale à celle du groupe de Picard.

\subsection*{2. Schémas $\underline{\operatorname{Hom}}_S(X,Y)$,
$\Pi_{X/S}Z$, $\underline{\operatorname{Aut}}_S(X)$, etc.}
\addcontentsline{toc}{subsection}{2. Schémas Hom, Pi et Aut}

Soient $X,Y$ deux pré-schémas sur $S$ ; pour tout pré-schéma $T$ sur $S$,

\src{381}
soient $X_T=X\times_ST$, $Y_T=Y\times_ST$, et considérons l'ensemble
\[
  F(T)=\operatorname{Hom}_T(X_T,Y_T)
      =\operatorname{Hom}_S(X_T,Y)
      =\operatorname{Hom}_S(X\times_ST,Y)
\]
comme un foncteur contravariant en $T$. S'il est représentable, on dénote par
$\underline{\operatorname{Hom}}_S(X,Y)$ le pré-schéma sur $S$ qui le
représente, donc on a un isomorphisme fonctoriel
\[
  \operatorname{Hom}_S
  \bigl(T,\underline{\operatorname{Hom}}_S(X,Y)\bigr)
  \xrightarrow{\sim}
  \operatorname{Hom}_S(T\times_SX,Y).
\]
Il y a bien des variantes de ce problème universel, dont la solution s'y
ramène : pré-schéma des $S$-automorphismes d'un $S$-pré-schéma $X$ (ce sera un
pré-schéma en groupes), pré-schéma des $S$-homomorphismes d'un $S$-pré-schéma
en groupes dans un autre (ce sera un pré-schéma en groupes commutatifs si le
deuxième schéma en groupes est commutatif), etc. On peut aussi généraliser la
définition de $\underline{\operatorname{Hom}}_S(X,Y)$ en considérant un
pré-schéma $Z$ sur le pré-schéma $X$ sur $S$, et le foncteur
\[
  F(T)=\operatorname{Hom}_{X_T}(X_T,Z_T)
\]
(ensemble des « sections » de $Z_T$ sur $X_T$) ; si ce foncteur est
représentable, le $S$-pré-schéma qui le représente sera noté $\Pi_{X/S}Z$, on
aura donc par définition un isomorphisme fonctoriel
\[
  \operatorname{Hom}_S(T,\Pi_{X/S}Z)
  =\operatorname{Hom}_{X_T}(X_T,Z_T).
\]
Faisant $Z=Y\times_SX$, on retrouve
$\underline{\operatorname{Hom}}_S(X,Y)$. De ces définitions résulte pour les
nouveaux pré-schémas ainsi introduits un formulaire aussi trivial qu'utile,
que nous ne donnerons pas ici (vu qu'il vaut dans toute catégorie où les
produits et produits fibrés existent). Plus sérieuse est la question de
l'existence des schémas du type $\underline{\operatorname{Hom}}_S(X,Y)$. On
constate d'abord que, pour $X$ fixé,
$\underline{\operatorname{Hom}}_S(X,Y)$ ne peut exister pour tout $Y$ sur $S$
(ou $\Pi_{X/S}Z$ pour tout $Z$ sur $X$) que si $X$ est plat sur $S$. De plus,
on se convainc qu'il n'est raisonnable de s'attendre à l'existence d'une
solution, pour des $Y$ assez généraux, que si $X$ est de plus propre sur $S$.
Il semble par ailleurs que ces conditions soient suffisantes pour l'existence
de $\underline{\operatorname{Hom}}_S(X,Y)$ et $\Pi_{X/S}Z$, à condition, le
cas échéant, de faire une hypothèse genre « quasi projective » sur $Y/S$, resp.
sur $Z/X$ ; c'est ce qu'on peut vérifier en tous cas dans des cas assez
nombreux (par exemple, lorsque $Y$ est affine sur $S$ ; ou, par constructions
directes élémentaires, lorsque $X$ est fini sur $S$). Voici ce que donnent les
théorèmes 1 et 2 :

\src{382}
\result{PROPOSITION}{2.1} Soient $\Lambda$ un anneau noethérien, $X$ et $Y$
deux pré-schémas quelconques sur $\Lambda$, considérons le foncteur
\[
  F(A)=\operatorname{Hom}_A(X_A,Y_A)
\]
sur la catégorie $\mathcal C_0$ des $\Lambda$-algèbres artiniennes. Si $X$ est
plat sur $\Lambda$, ce foncteur est pro-représentable.

De plus, on vérifie que pour tout $A\in\mathcal C_0$ et $\xi\in F(A)$, on a un
isomorphisme canonique
\[
  F(I_A,\xi)=H^0\!\left(
    X_A,
    \underline{\operatorname{Hom}}_{\mathcal O_{X_A}}
    \bigl(\xi^*(\Omega^1_{Y_A/A}),\mathcal O_{X_A}\bigr)
  \right),
\]
où $\Omega^1_{Y_A/A}$ est le faisceau des $1$-différentielles de Kähler de
$Y_A$ par rapport à $A$. Prenant pour $A$ un corps, on trouve, utilisant
A.5.1 et le théorème de finitude de [2], le corollaire suivant :

\result{COROLLAIRE}{} Supposons $X$ plat et propre sur $S$, $Y$ localement de
type fini sur $S$. Alors $F$ est pro-représentable et les composants locaux de
la $\Lambda$-algèbre topologique correspondante sont des anneaux noethériens.

\result{REMARQUES}{} Les problèmes envisagés dans ce numéro, et de nombreux
autres, étaient étudiés communément, dans le cadre de la géométrie algébrique
classique, à l'aide des « coordonnées de Chow » des cycles dans l'espace
projectif, permettant de considérer ces cycles comme des points de variétés
projectives convenables. Ce procédé, et de façon générale l'utilisation des
coordonnées de Chow, semble irrémédiablement insuffisant dans le point de vue
des schémas, car il détruit les éléments nilpotents dans les variétés de
paramètres, et en particulier ne se prête pas à une étude satisfaisante des
variations infinitésimales de cycles (sans compter sa nature non intrinsèque,
liée à l'espace projectif). Le langage des coordonnées de Chow était
malheureusement le seul utilisé par de nombreux géomètres algébristes pour
l'étude des familles de variétés ou des familles de cycles, ce qui semble
avoir été un obstacle sérieux à la clarification de ces notions, malgré son
intérêt technique certain (probablement provisoire). Si on veut obtenir
l'analogue des variétés de Chow en théorie des schémas, on est conduit au
problème universel suivant : soit $X$ un pré-schéma sur $S$, pour tout
pré-schéma $T$ sur $S$, on considère l'ensemble $F(T)$ des sous-pré-schémas
fermés de $X_T=X\times_ST$ qui sont plats sur $T$. On essaie de représenter ce
foncteur en $T$ à l'aide d'un pré-schéma sur $S$. Plus généralement, on peut se
donner un faisceau quasi cohérent $G$ sur $X$, et prendre pour $F(T)$
l'ensemble des faisceaux

\src{383}
quotients de $G_T$ qui sont plats sur $T$. Il semble qu'il existe une solution
du problème, avec un schéma $C$ localement de type fini sur $S$, lorsque $X$
est propre sur $S$ localement noethérien et lorsque $G$ est de plus cohérent.
En tous cas, supposant seulement $S$ localement noethérien, la restriction de
$F$ aux « $S$-algèbres artiniennes » est pro-représentable, et si de plus $X$
est propre sur $S$ et $G$ cohérent, les composants locaux de l'anneau
topologique $O$ correspondant sont noethériens. Bien entendu, même une fois
prouvée l'existence du « schéma de Chow » de $X$ sur $S$, il restera à en
trouver une décomposition en ouverts disjoints $C_i$ (correspondants à la
fixation d'invariants continus tels que degré et dimension des cycles qu'on
fait varier), qui soient de type fini (et non seulement localement de type
fini) sur $S$, à préciser les relations de ce schéma avec les classiques
variétés de Chow et à préciser quand un $C_i$ est même projectif ou du moins
quasi projectif sur $S$.\footnote{L'erratum de 1962 précise que le cas
projectif est entièrement résolu par les schémas de Hilbert (TDTE IV), tandis
que des exemples de Nagata et Hironaka montrent que, sans hypothèse
projective, ces foncteurs peuvent ne pas être représentables, même pour les
sous-variétés de dimension zéro d'une variété complète non singulière de
dimension trois ; ce phénomène est lié à la possible inexistence du carré
symétrique.}

\subsection*{3. Schémas de Picard}
\addcontentsline{toc}{subsection}{3. Schémas de Picard}

Soit $f:X\longrightarrow S$ un $S$-pré-schéma,\footnote{Pour une étude plus
complète, l'erratum de 1962 renvoie à TDTE V.} considérons le faisceau
multiplicatif $\mathcal O_X^*$ des unités du faisceau structural de $X$, et le
groupe
\[
  P(X/S)=H^0\bigl(S,R^1f_*(\mathcal O_X^*)\bigr),
\]
appelé groupe de Picard relatif de $X/S$. Un élément de ce groupe est donc
défini en se donnant un recouvrement ouvert $(U_i)$ de $S$, et un faisceau
inversible $L_i$ sur chaque $f^{-1}(U_i)$, de telle façon que, pour tout couple
$(i,j)$, $L_i|f^{-1}(U_i\cap U_j)$ soit isomorphe à
$L_j|f^{-1}(U_i\cap U_j)$, du moins localement au-dessus de $U_i\cap U_j$
(i.e. ces deux faisceaux sont « équivalents » au sens de [3], B, 4). Si $X/S$
admet une section, alors $P(X/S)$ n'est autre que l'ensemble des classes de
faisceaux inversibles sur $X/S$ à « équivalence » près (loco citato). Posons
maintenant, pour tout $T$ sur $S$,
\[
  F(T)=P(X_T/T) ;
\]
on obtient un foncteur contravariant en $T$, qu'on pourra appeler le foncteur de
Picard de $X/S$ ; si ce foncteur est représentable, le pré-schéma sur $S$ qui
le représente sera appelé le pré-schéma de Picard de $X/S$ et noté
$\underline P(X/S)$. On aura donc un isomorphisme de foncteurs
\[
  \operatorname{Hom}_S\bigl(T,\underline P(X/S)\bigr)
  \xrightarrow{\sim}P(X_T/T).
\]
La formation de pré-schémas de Picard est compatible avec l'extension de la
base, en particulier les pré-schémas de Picard des fibres de $X$ sur $S$ (qui
sont des pré-schémas sur les corps résiduels $K(s)$ des $s\in S$) sont les
fibres de $\underline P(X/S)$. Bien entendu, comme $P(X_T/T)=F(T)$ est un
groupe commutatif, les pré-schémas de Picard sont des pré-schémas en groupes.
Notons aussi que les jacobiennes généralisées de Rosenlicht ne sont autres que
les composantes connexes de l'unité dans les schémas de Picard de courbes
complètes pouvant avoir des singularités, ce qui devrait rendre évidentes la
plupart de leurs propriétés (une fois démontrée l'existence).

\result{REMARQUE}{} La définition adoptée ici n'est raisonnable que lorsque
tout point de $S$\footnote{Le texte imprimé porte « $Y$ » ; aucune base $Y$
n'est définie ici, et la condition porte sur la base $S$.} admet un voisinage
ouvert $U$ sur lequel $X$ admette une section. Dans le cas général, il faut
modifier un peu la définition du foncteur de Picard, pour pouvoir obtenir
encore un théorème d'existence.

Ici, les conditions plausibles d'existence d'un pré-schéma de Picard sont les
suivantes : $X$ est propre et plat sur $S$, $f_*(\mathcal O_X)=\mathcal O_S$,
et $X$ admet localement une section sur $S$. Cette condition s'introduit de
façon naturelle dans l'application de la technique de descente, pour éliminer
les automorphismes d'un faisceau inversible $L$ sur $X$ en le munissant d'une
section marquée au-dessus de la section $s$ ([3], B, 4). On trouve notamment :

\src{384}
\result{PROPOSITION}{3.1} Supposons $X$ plat sur
$S=\operatorname{Spec}(\Lambda)$, $\Lambda$ noethérien, et que, pour tout $T$
de type fini sur $S$, on ait
\[
  f_{T*}(\mathcal O_{X_T})=\mathcal O_T
\]
(si $f$ est propre et séparable,\footnote{L'erratum de 1962 prescrit
explicitement l'hypothèse « $f$ propre et séparable ».} ou si $S$ est le
spectre d'un corps, il résulte de Künneth que la dernière condition équivaut à
$f_*(\mathcal O_X)=\mathcal O_S$). Alors le foncteur de Picard de $X/S$ sur la
catégorie des $\Lambda$-algèbres artiniennes est pro-représentable.

De plus, on aura ici
\[
  F(I_A,\xi)=H^1(X_A,\mathcal O_{X_A}),
\]
en particulier :

\result{COROLLAIRE}{} Si $X$ est propre sur $S$, alors les composants locaux
de la $\Lambda$-algèbre topologique $O$ correspondant au foncteur de Picard
sont noethériens.

\result{REMARQUES}{} On peut généraliser les définitions et résultats de ce
numéro à la classification des fibrés principaux sur $X$, à groupe structural
un schéma en groupes $G$ sur $S$ qui est affine et plat sur $S$, et
commutatif. Dans le cas où $G$ ne serait pas commutatif, donc le fibré en
groupes adjoint d'un fibré principal (dont les sections sont les
automorphismes du fibré principal) n'est plus trivial, la proposition 3.1 ne
reste plus valable telle quelle. On peut cependant modifier le problème
universel de façon à obtenir encore une solution (du moins, pour l'instant, en
géométrie formelle). La règle d'or à retenir, dans le contexte du présent
\src{385}numéro et des suivants, et chaque fois qu'on cherche des « schémas de modules »
pour des classes d'objets qui ne sont définis qu'à un isomorphisme près, reste
toujours la suivante : éliminer les automorphismes éventuels des objets qu'on
veut classifier, par l'introduction, si besoin est, de structures auxiliaires
(points ou éléments de sections marquées, fixation de formes différentielles,
etc.), qu'on prendra assez anodines pour ne pas modifier de façon substantielle
le problème initial.

\subsection*{4. Modules formels d'une variété}
\addcontentsline{toc}{subsection}{4. Modules formels d'une variété}

Soit $\Lambda$ un anneau local noethérien de corps résiduel $k$ (le plus
souvent, $\Lambda$ sera égal à $k$, ou à un $p$-anneau de Cohen), $X_0$ un
pré-schéma sur $k$. Pour toute $\Lambda$-algèbre artinienne locale $A$, on
considère l'ensemble $F(A)$ des classes, à un isomorphisme près, de
$A$-pré-schémas $X$ plats sur $A$, munis d'un isomorphisme
\[
  (*)\qquad
  X\otimes_A\kappa(A)\xleftarrow{\sim}
  X_0\otimes_k\kappa(A),
\]
où $\kappa(A)$ est le corps résiduel de $A$ ; bien entendu, les isomorphismes
entre tels $A$-pré-schémas plats doivent respecter l'isomorphisme précédent
donné dans la structure. Si $A$ est une $\Lambda$-algèbre artinienne non
nécessairement locale, de composantes locales $A_i$, on prend pour $F(A)$ le
produit des $F(A_i)$. Ainsi $F$ devient un foncteur multiplicatif en $A$, qu'on
pourra appeler le foncteur des modules pour $X_0$ (et $\Lambda$). Si ce
foncteur est représentable, il lui correspond une $\Lambda$-algèbre
topologique $O$ locale, de corps résiduel $k$, et le spectre formel de $O$ sera
appelé schéma formel des modules pour $X_0$ (et $\Lambda$), cf. [2] pour
quelques détails sur cette notion.

Ici, pour appliquer la technique de descente, les automorphismes « finis » de
$X_0$ sont inoffensifs, car ils sont sans influence sur l'existence
d'automorphismes (au sens précisé plus haut) des $A$-pré-schémas $X$ ; la
condition nécessaire et suffisante, lorsque $A$ n'est pas réduit à un corps,
pour que $X$ n'ait pas de $A$-automorphisme non trivial, est qu'on ait
\[
  H^0(X_0,\Theta_{X_0/k})=0,
\]
où $\Theta_{X_0/k}$ désigne le faisceau des $k$-dérivations (= faisceau
tangent) de $X_0$. On trouve d'ailleurs facilement (du moins si $X_0$ est
simple sur $k$)
\[
  F(I_A,\xi)=H^1(X_A,\Theta_{X_A/A}).
\]
On en conclut alors, comme d'habitude :

\src{386}
\result{PROPOSITION}{4.1} Supposons que
$H^0(X_0,\Theta_{X_0/k})=0$. Alors le schéma formel des modules pour $X_0$
existe. Si de plus $X_0$ est propre sur $k$, le schéma formel des modules est
noethérien.

\result{REMARQUES}{} 1\textsuperscript{o} Lorsque $X_0$ n'est pas supposé
simple sur $k$, $F(I_A,\xi)$ s'identifie à un sous-$A$-module de
\[
  \operatorname{Ext}^1_{\mathcal O_{P_A}}
  \bigl(P_A;\mathcal I_{X_A},\mathcal O_{X_A}\bigr),
\]
où on pose $P_A=X_A\times_\Lambda X_A$, où $\mathcal O_{X_A}$ est regardé
comme un faisceau cohérent sur $P_A$ grâce au morphisme diagonal
$X_A\longrightarrow P_A$, et où $\mathcal I_{X_A}$ désigne le faisceau
cohérent d'idéaux sur $P_A$ défini par le morphisme diagonal. De façon précise,
une globalisation facile de la théorie de Hochschild montre que le
$\operatorname{Ext}^1$ ci-dessus s'identifie à l'ensemble des classes, à un
isomorphisme près, de faisceaux de $I_A$-algèbres plates $\mathcal O$ sur
$X_A$, munies d'un isomorphisme d'augmentation
\[
  \mathcal O\otimes_{I_A}A\longrightarrow\mathcal O_{X_A}
\]
(on rappelle que l'on pose $I_A=A[t]/(t^2)$). Le sous-module $F(I_A,\xi)$ est
celui qui correspond aux faisceaux d'algèbres commutatives. Les hypothèses de
simplicité ne sont donc pas essentielles en théorie des modules, comme [2] le
laissait entendre.

2\textsuperscript{o} Rappelons (loco citato) qu'en particulier toute courbe
algébrique $X_0$ simple et propre sur $k$ admet un schéma formel des modules,
simple sur $\Lambda$, de dimension relative $3g-3$ si le genre $g$ est
$\geq2$, $g$ si $g=0$ ou $1$. Ces deux derniers cas ne rentrent plus
directement dans la proposition 4.1. On pourrait cependant s'y ramener, dans
le cas des courbes elliptiques ($g=1$), grâce aux considérations qui vont
suivre.

On peut, bien entendu, varier ad libitum la proposition 4.1 en considérant des
systèmes de schémas sur $k$, munis de structures diverses. Supposons, par
exemple, que $X_0$ soit un schéma abélien sur $k$, à point origine marqué
(i.e. $X_0$ est considéré comme un schéma en groupes sur $k$), et soit $F(A)$
l'ensemble des classes, à un isomorphisme près, de schémas abéliens sur $A$
(i.e. de schémas en groupes propres et simples sur $A$) munis d'un
isomorphisme de schémas abéliens $(*)$. On vérifie que l'imposition d'une
structure multiplicative (et même seulement d'une « section unité ») élimine
les automorphismes infinitésimaux, et qu'il existe par suite un schéma formel
de modules, correspondant à un anneau local complet noethérien $O$. On peut
d'ailleurs montrer que, si $X$ est un schéma propre et simple à fibres
« absolument connexes » au-dessus d'un pré-schéma localement noethérien $S$,
toute structure multiplicative sur $X$ admettant une section unité est
nécessairement associative et commutative (pourvu du moins qu'elle le soit sur
une fibre et que $S$ soit connexe) et est de plus uniquement déterminée par la
connaissance

\src{387}
de la section unité. De plus, supposant que $S$ est le spectre d'un anneau
local artinien $A$ de corps résiduel $k$, $X$ propre sur $A$ et muni d'une
section $s$, et enfin $X\otimes_Ak$ muni d'une structure de schéma abélien sur
$k$, admettant pour élément neutre le point de $X\otimes_Ak$ correspondant à
$s$, un calcul facile d'obstructions joint à un raisonnement dû à Serre permet
de prouver qu'il existe effectivement sur $X$ une structure multiplicative
admettant la section $s$ comme section unité. (À partir de là, utilisant le
« théorème d'existence » de [2] pour passer au cas où $A$ est noethérien local
complet, puis la technique de descente de [3] pour le cas général, on peut
prouver l'énoncé analogue pour tout $S$ localement noethérien connexe.) Cela
prouve que le foncteur $F(A)$ considéré ici est isomorphe au foncteur analogue
défini au début de ce numéro en faisant abstraction de la structure
multiplicative sur $X_0$. Il s'ensuit en particulier que, si $\mathfrak m$ est
l'idéal maximal de $O$, $\mathfrak m/\mathfrak m^2$ est canoniquement
isomorphe au dual de $H^1(X_0,\Theta_{X_0/k})$ et est donc de dimension $n^2$,
où $n=\dim X_0$. Il serait fort intéressant de déterminer si $O$ est bien
simple sur $\Lambda$, i.e. isomorphe à une algèbre de séries formelles à $n^2$
indéterminées sur $\Lambda$. Le paragraphe 1, proposition 5.2, permet de donner
une formulation équivalente de ce problème comme un problème d'existence de
schémas abéliens se réduisant suivant un schéma abélien donné. On voit en tous
cas, par voie transcendante, que la réponse est affirmative si $k$ est de
caractéristique zéro. En caractéristique $p\ne0$, il suffit évidemment de se
borner au cas où $\Lambda$ est l'anneau des vecteurs de Witt construit sur un
corps algébriquement clos $k$. Ce serait peut-être le moment pour le « foncteur
de Greenberg » de faire ses preuves\ldots\footnote{L'erratum de 1962 annonce
que le schéma formel des modules d'une variété abélienne sur un corps est
simple sur l'anneau de Witt : tout schéma abélien sur un quotient artinien
local se relève à l'anneau dont il est quotient. La démonstration utilise la
variance de la classe d'obstruction au relèvement introduite dans l'exposé 182,
p. 12. Il rappelle aussi les constructions de Mumford pour les courbes de
genre $g$ et les schémas abéliens polarisés (cf. SHMT).}

\subsection*{5. Extension des revêtements}
\addcontentsline{toc}{subsection}{5. Extension des revêtements}

Soient $\mathfrak X$ un pré-schéma formel noethérien [2], $U$ une partie
ouverte de $\mathfrak X$, définie localement par la « non-annulation » d'une
section de $\mathcal O_{\mathfrak X}$ qui est non-diviseur de zéro non
seulement sur $\mathfrak X$, mais sur tout
$X_n$,\footnote{Précision imposée par l'erratum de 1962.} donc assez gros pour
que toute section de $\mathcal O_{\mathfrak X}$ sur un ouvert $V$, nulle sur
$U\cap V$, soit nulle. Soit $\mathcal J$ un « idéal de définition » pour
$\mathfrak X$, et soit
\[
  X_n=(\mathfrak X,\mathcal O_{\mathfrak X}/\mathcal J^{n+1}),
\]
qui est donc un pré-schéma noethérien ordinaire. Si alors
$\mathfrak X',\mathfrak X''$ sont deux revêtements plats de $\mathfrak X$
(i.e. deux pré-schémas sur $\mathfrak X$ définis par des faisceaux d'algèbres
qui sont cohérents et localement libres comme faisceaux de modules), non
ramifiés sur $U$, l'application évidente
\[
  \operatorname{Hom}_{\mathfrak X}(\mathfrak X',\mathfrak X'')
  \longrightarrow
  \operatorname{Hom}_{X_0}(X'_0,X''_0)
\]
est injective ; en particulier, un automorphisme de $\mathfrak X'$ qui induit
l'identité sur $X'_0$ est l'identité. Cela permet d'appliquer à la situation la
technique de descente. Partons, en particulier, d'un revêtement plat $X'_0$ de
$X_0$, non ramifié

\src{388}
au-dessus de $U_0$, soit $G(\mathfrak X)$ l'ensemble des classes, à un
isomorphisme près (induisant l'identité sur $X'_0$), de revêtements plats
$\mathfrak X'$ de $\mathfrak X$ qui induisent $X'_0$ sur $X_0$ (et sont donc
nécessairement non ramifiés sur $U$). On définit de même $G(V)$ pour toute
partie ouverte $V$ de $\mathfrak X$, et plus généralement $G(\mathfrak Y)$
pour tout pré-schéma formel $\mathfrak Y$ au-dessus de $\mathfrak X$. Cela dit,
les résultats de [2] et [3] impliquent d'abord les résultats suivants :

\noindent a. Lorsque $V$ est un ouvert variable de $\mathfrak X$, les $G(V)$
forment un faisceau sur $\mathfrak X$, soit $G_{\mathfrak X}=G$. La
restriction de ce faisceau à $U$ est le faisceau constant dont les fibres sont
réduites à un élément.

De façon plus générale, la détermination des fibres de $G_{\mathfrak X}$ est
une question d'anneaux locaux complets, de façon précise :

\noindent b. Pour tout $x\in\mathfrak X$, on a
\[
  G_x=G\bigl(\operatorname{Spec}(\mathcal O_{\mathfrak X,x})\bigr)
  \subset
  G\bigl(\operatorname{Spec}(\widehat{\mathcal O}_{\mathfrak X,x})\bigr),
\]
(= classes à un isomorphisme près d'algèbres finies et libres $B$ sur
$\widehat{\mathcal O}_{\mathfrak X,x}$ munies d'un isomorphisme de
$B\otimes_{\widehat{\mathcal O}_{\mathfrak X,x}}\mathcal O_{X_0,x}$ avec
$\widehat{\mathcal O}'_{0x}$, où $\mathcal O'_0$ est le faisceau d'algèbres sur
$X_0$ qui définit $X'_0$).\footnote{La formule imprimée identifie directement
$G_x$ au terme complété et omet le chapeau sur $\mathcal O'_{0x}$. L'inclusion
et le complété sont ceux prescrits par l'erratum de 1962.}

\noindent c. On a un isomorphisme canonique
\[
  G_{\mathfrak X}=\varprojlim_nG_{X_n} ;
\]
en d'autres termes, pour tout ouvert $V$ de $\mathfrak X$, on a
\[
  G(V)=\varprojlim_nG(V_n).
\]

\noindent d. Supposons que $\mathfrak X$ se déduise d'un schéma ordinaire
propre $X$ sur un anneau local noethérien complet $\Lambda$ ayant un idéal de
définition $\mathfrak m$, en prenant le complété $\mathcal J$-adique de
$\mathcal O_X$, où $\mathcal J=\mathfrak m\mathcal O_X$. Alors
$G(\mathfrak X)$ est canoniquement isomorphe à l'ensemble des classes de
revêtements plats du schéma ordinaire $X$ qui « se réduisent suivant $X'_0$ ».

De façon imagée, on peut dire que (a) et (b) établissent les relations
fondamentales entre l'aspect local et l'aspect global du problème, (c) donne
les relations entre l'aspect « fini » et l'aspect « infinitésimal », enfin (d)
rappelle (sous les conditions précisées) l'identité entre l'aspect « formel »
et l'aspect « algébrique ».

Supposons maintenant que $\mathfrak X$ soit défini sur un anneau local
noethérien complet $\Lambda$, avec
$\mathcal J=\mathfrak m\mathcal O_{\mathfrak X}$ (donc $X_0$ est un
pré-schéma sur $\Lambda/\mathfrak m$). Pour toute algèbre $A$ finie sur
$\Lambda$, posons
\[
  F(A)=G(\mathfrak X\times_\Lambda A),
\]
où le marquage est pris après changement de base à l'algèbre réduite $A_{\mathrm{red}}$, et non littéralement à $A/\mathfrak mA$. C'est un foncteur covariant en $A$, à valeurs dans la catégorie des ensembles,
et d'après (c) ce foncteur est complètement connu si on le connaît pour $A$
artinien ; et il revient au même de dire que ce foncteur est pro-représentable,
i.e. de la forme
\[
  F(A)=\operatorname{Hom}_{\Lambda\text{-algèbres top.}}(O,A),
\]

\src{389}
où $O$ est une $\Lambda$-algèbre topologique du type envisagé dans A,
paragraphe 5, ou de dire qu'il en est de même quand on se restreint aux $A$
artiniens. La conjonction des théorèmes 1 et 2 implique alors effectivement :

\result{PROPOSITION}{5.1} Pourvu que toute famille artinienne classifiée par le foncteur soit plate sur son algèbre de coefficients et que l'ouvert choisi reste schématiquement dense après les changements de base pertinents, le foncteur précédent est pro-représentable.

Bien entendu, d'après (a), si $U=\mathfrak X$, $G(\mathfrak Y)$ est réduit à
un seul élément pour tout $\mathfrak Y$ sur $\mathfrak X$, et le foncteur $F$
est alors sans grand intérêt (on aura $O=\Lambda$). Il semble que, dans
pratiquement tout autre cas, l'anneau local topologique $O$ n'est pas
noethérien. Son existence met néanmoins en évidence, de façon frappante, la
nature « continue » de l'ensemble $G(\mathfrak X)$ des solutions
(correspondant intuitivement au fait qu'il y a un choix « continu » dans la
façon dont la ramification se propage quand on fait une extension de $X'_0$).
On comparera ce résultat avec le point de vue de J.-P. Serre [4] en théorie du
corps de classes local, mettant lui aussi en évidence le caractère continu du
groupe de Galois topologique de l'extension abélienne maximale d'un corps local
« géométrique », le groupe dual (au sens de Pontrjagin) y apparaissant comme
une limite inductive de groupes algébriques (ou du moins quasi algébriques) ;
ici aussi, la classification des extensions est donnée par des « variétés » de
dimension infinie. D'ailleurs, on peut prendre dans ce qui précède pour
$\mathfrak X$ le spectre formel d'un anneau local complet (dont $\Lambda$ sera,
par exemple, un sous-anneau de Cohen), et on peut espérer que les résultats de
ce numéro puissent être utilisés dans l'étude des extensions d'un anneau local
complet de dimension $>1$. Aussi bien dans le cas local que global, ils
permettront peut-être de formuler des relations précises entre les phénomènes
de ramification supérieure et des phénomènes en caractéristique zéro
(abordables par voie transcendante). En tous cas, c'est l'analyse préliminaire
à la proposition 5.1 qui permet d'étendre au cas « tamely ramified » les
méthodes exposées dans [2] pour l'étude du groupe fondamental, et de résoudre
par voie transcendante le « problème des trois points ».

Pour finir, signalons que la situation se simplifie lorsque $X_0$ est de dimension 1 : alors, en vertu de (a), complété par le recollement formel sur les schémas locaux complétés épointés,
\[
  G(\mathfrak X)\simeq
  \prod_iG\bigl(\operatorname{Spec}(\widehat{\mathcal O}_{\mathfrak X,x_i})\bigr),
\]
où les $x_i$ sont les points de $X_0-U$ : on peut se donner arbitrairement les
extensions « locales » aux points de ramification. La normalité de $X_0$ ne suffit pas à rendre simple sur $\operatorname{Spec}(\Lambda)$ le schéma formel des modules de 5.1 ; la dernière assertion imprimée est fausse en général.

\src{390}
\begin{center}
  \textbf{BIBLIOGRAPHIE}
\end{center}

\begin{description}
  \item[{[1]}] Grothendieck (Alexander). --- \emph{Sur quelques points
    d'algèbre homologique}, Tohoku Math. J., t. 9, 1957, p. 119--221.
  \item[{[2]}] Grothendieck (Alexander). --- \emph{Géométrie formelle et
    géométrie algébrique}, Séminaire Bourbaki, t. 11, 1958/59, fasc. 3,
    n\textsuperscript{o} 182.
  \item[{[3]}] Grothendieck (Alexander). --- \emph{Technique de descente et
    théorèmes d'existence en géométrie algébrique, I : Généralités, Descente
    par morphismes fidèlement plats}, Séminaire Bourbaki, t. 12, 1959/60,
    fasc. 1, n\textsuperscript{o} 190.
  \item[{[4]}] Serre (Jean-Pierre). --- \emph{Corps locaux et isogénies},
    Séminaire Bourbaki, t. 11, 1958/59, fasc. 3, n\textsuperscript{o} 185.
\end{description}

\par\smallskip
\begingroup\footnotesize
\endgroup
